% Paper B: Accessibility, Stability, and Emergent Geometry
% Conceptual Clarifications on the Holographic Alaya-Field Framework
% FINAL REVISED VERSION

\documentclass[12pt,a4paper]{article}
\usepackage{arxiv}
\usepackage{amsmath,amssymb,amsfonts,amsthm}
\usepackage{physics}
\usepackage{hyperref}
\usepackage{tcolorbox}
\usepackage{array}
\newcolumntype{P}[1]{>{\raggedright\arraybackslash}p{#1}}

\newtheorem{theorem}{Theorem}[section]
\newtheorem{definition}[theorem]{Definition}
\newtheorem{remark}[theorem]{Remark}

\title{Accessibility, Stability, and Emergent Geometry:\\
Conceptual Clarifications on the Holographic Alaya--Field Framework}

\author{
  Sidong Liu, PhD \\
  iBioStratix Ltd \\
  \texttt{sidongliu@hotmail.com}
}

\date{July 2026}

\begin{document}
\emergencystretch=2em
\raggedbottom
\hbadness=5000
\vbadness=5000

\maketitle

\begin{abstract}
This paper provides conceptual clarification of the Holographic Alaya-Field Framework (HAFF) introduced in our previous work. 
We address three potential misreadings: subjectivism (that observers create spacetime), anti-realism (that geometry is illusory), and trivialism (that the framework reduces to coordinate choice). 
By analyzing the structural notion of accessibility via stability conditions, delineating boundaries with existing interpretations (AQFT, RQM, QBism, MWI), and characterizing emergent geometry as a stable organizational phase, we clarify what the framework commits to and what it does not. 
This analysis is purely interpretational; no new formal results are introduced.
\end{abstract}

\section{Introduction}
\label{sec:intro}

Our previous work proposed a framework in which effective operator and connectivity descriptions depend on a selected algebra acting on a global quantum state \cite{Liu2026}.  Metric inequivalence is conditional on a reconstruction dictionary; it is not a general consequence of channel or algebra inequivalence.
However, structural novelty of this kind is particularly vulnerable to interpretational confusion. 
The present paper addresses these interpretational implications and clarifies the conceptual commitments of the framework.

\subsection{The Risk of Misreading}

The claim that geometry is coarse-graining-dependent naturally invites several misreadings, each of which conflates distinct notions of dependence. 
Three such misreadings are especially common:

\begin{enumerate}
\item \emph{Subjectivism}: The view that observers or agents create spacetime through their choices or beliefs, collapsing the framework into an epistemic or observer-relative interpretation.
\item \emph{Anti-realism}: The view that spacetime has no objective existence, and that emergent geometry is therefore illusory or merely pragmatic.
\item \emph{Trivialism}: The view that coarse-graining-dependent geometry reduces to a choice of coordinates or descriptive convention, with no substantive physical consequences.
\end{enumerate}

Each of these readings is incorrect, but each arises naturally from surface-level features of the formalism. 
The purpose of this paper is to block such misreadings by clarifying the nature of accessibility, the structural role of observable algebras, and the ontological status of emergent geometry within the framework.

\subsection{Scope and Objectives}

This paper introduces no theorems.  The stability criteria and phase quantities below are explicitly heuristic diagnostics and model inputs, not proved selection principles.  The paper provides a systematic interpretational analysis aimed at three specific objectives:

\begin{enumerate}
\item \emph{Clarify the notion of accessibility}: We formulate accessibility as a structural and operational modelling proposal relative to declared dynamics, probes, and mereology, rather than as an epistemic or observer-dependent notion.
\item \emph{Delineate boundaries with existing interpretations}: We situate the framework in relation to algebraic quantum field theory, relational quantum mechanics, QBism, and the Many-Worlds interpretation, clarifying points of agreement, divergence, and complementarity.
\item \emph{Characterize the ontological status of emergent geometry}: We argue that geometry functions as a stable organizational phase of quantum information, analogous to phases in condensed matter systems, avoiding both naive realism and anti-realist eliminativism.
\end{enumerate}

Importantly, this analysis does not constitute a defense of the framework, nor does it aim to persuade readers of its correctness. 
The goal is clarity: to ensure that the structural commitments of the framework are understood on their own terms, and that criticisms, if any, are directed at what the framework actually claims rather than at interpretational projections.

\subsection{What This Paper Does Not Do}

To further delimit scope, we note explicitly what this paper does \emph{not} attempt:

\begin{itemize}
\item It does not propose new dynamics, empirical predictions, or modifications to quantum mechanics.
\item It does not claim that the framework resolves outstanding problems in quantum gravity, quantum foundations, or the measurement problem.
\item It does not advocate for any particular metaphysical or philosophical position beyond the minimal structural commitments required by the formalism itself.
\item It does not interpret the framework as implying idealism, observer-created reality, or any form of mind-dependence.
\end{itemize}

The analysis remains strictly within the domain of structural interpretation: identifying what the mathematical formalism commits to, what it leaves open, and how it relates to existing approaches.

\subsection{Organization}

The paper proceeds as follows. Section~\ref{sec:stance} briefly recapitulates the structural stance adopted in our previous work, emphasizing the priority of observable algebras over tensor factorizations and the role of coarse-graining in defining effective subsystems. 

Section~\ref{sec:accessibility} provides a detailed analysis of candidate accessibility diagnostics and their physical inputs. A taxonomy distinguishes subjective, relational, and structural notions of dependence; HAFF adopts the third as a model-relative proposal without claiming a unique input-free selector.

Section~\ref{sec:relations} examines the relationship between the framework and four representative approaches: algebraic quantum field theory, relational quantum mechanics, QBism, and the Many-Worlds interpretation. Each subsection clarifies points of conceptual overlap and structural divergence, preventing conflation while identifying opportunities for complementarity.

Section~\ref{sec:geometry_phase} argues that emergent geometry should be understood as a stable organizational phase of entanglement structure, drawing on analogies with condensed matter physics. This perspective avoids treating geometry as either fundamental or illusory, instead characterizing it as contingent but objective—dependent on physical conditions rather than epistemic contexts.

Section~\ref{sec:scope} delimits the structural assumptions of the framework, enumerates what it does not commit to, and outlines open questions for future investigation. Particular attention is given to the distinction between structural analysis and metaphysical advocacy.

\subsection{Methodological Note}

Throughout this paper, we adopt a deliberately conservative rhetorical stance. Claims are hedged with modal qualifiers ("may suggest," "is consistent with," "can be understood as") because the selection criteria and phase analogy are heuristic and because interpretational conclusions should not outrun their model inputs. The goal is to present the framework as one coherent way of organizing the conceptual landscape, not as the uniquely correct interpretation.

This methodological caution reflects a broader commitment: interpretational clarity is best served by precision and restraint, not by advocacy or polemics. We aim to make the framework legible to researchers across different interpretational traditions, facilitating comparison and critique rather than preempting it.

\section{Structural Stance Recap}
\label{sec:stance}

We briefly recapitulate the structural stance of our previous work without repeating full mathematical derivations.

The framework uses three claims with different epistemic status:
\begin{enumerate}
\item \textbf{Factorization input}: Hilbert-space structure alone supplies no preferred tensor factorization.  Paper~A proves entanglement typicality only for a finite, physically preselected family; choosing that family remains a quantum-mereology input.

\item \textbf{Conditional channel separation}: Distinct finite-dimensional coarse-graining channels generically produce different outputs.  Entropy separation requires an explicit nondegeneracy condition.

\item \textbf{Conditional geometry}: Entropy/connectivity data become spacetime geometry only under a supplied reconstruction dictionary.  Non-diffeomorphism follows only when that dictionary separates the compared profiles modulo diffeomorphisms.
\end{enumerate}

The conceptual core can be summarized by reversing the standard explanatory arrow:

\begin{center}
\textbf{Standard:} State + Tensor factorization $\to$ Subsystems $\to$ Entanglement $\to$ Geometry
\end{center}

\begin{center}
\textbf{HAFF:} State + selected algebra/channel $\to$ Effective subsystems $\to$ Entanglement/connectivity $\xrightarrow{\text{dictionary}}$ Geometry
\end{center}

Write $\mathcal A_{\mathrm{sel},\mathbf c}$ for the selected algebra used in Paper~A.  It is not Paper~E's probe-reachable effect system, pointer/record algebra, or environmental fixed set.  The present paper investigates physical diagnostics that may constrain $\mathcal A_{\mathrm{sel},\mathbf c}$; it does not prove that the Hamiltonian uniquely determines it.

This reversal has a modest but consequential implication: subsystem structure, and consequently geometry, is not intrinsic to the quantum state alone but depends on which observables are accessible—where accessibility is understood in a precise, structural sense developed in the next section.

\section{Accessibility vs Observer-Dependence}
\label{sec:accessibility}

The notion of accessibility is central to the framework, and it is here that the risk of misreading is greatest.  HAFF intends accessibility diagnostics to concern physical stability and records rather than beliefs, but their outputs remain conditional on dynamics, probes, fragmentation, and candidate-algebra data.

\subsection{Three Notions of Dependence}

To prevent conflation, we distinguish three distinct senses in which a physical quantity might be said to "depend on" something:

\begin{table}[h]
\centering
\begin{tabular}{|l|l|l|l|}
\hline
\textbf{Type} & \textbf{Depends On} & \textbf{Example} & \textbf{Framework} \\
\hline
Subjective & Agent's beliefs/knowledge & Bayesian probability & QBism \\
Relational & Reference system & Velocity in SR & RQM \\
Structural & Physical interaction pattern & Decoherence basis & HAFF \\
\hline
\end{tabular}
\caption{Taxonomy of dependence notions}
\label{tab:dependence}
\end{table}

HAFF's notion of accessibility falls squarely in the third category.

\subsection{Heuristic Stability Diagnostics}

The following diagnostics constrain a supplied candidate algebra; they neither define measurement accessibility nor constitute sufficient conditions for $\operatorname{Fix}(\mathcal E)$.  Each diagnostic requires its own domain data.  In particular, the quantum-Darwinism diagnostic requires a system/environment fragmentation and the OTOC diagnostic requires declared probe operators.

Consider a unital candidate $\mathcal A\subset\mathcal B(\mathcal H)$ and a normal UCP Heisenberg dynamics $\mathcal E_t$.

\paragraph{Criterion 1: Dynamical Invariance.}
On a declared window $0\le t\le T$, require a scale-controlled bound
\begin{equation}
\sup_{0\le t\le T}
\sup_{O\in\operatorname{Ball}(\mathcal A)}
\|\mathcal E_t(O)-O\|\le\epsilon,
\end{equation}
or, where complete-bounded control is required, the corresponding CB-norm bound on $\mathcal E_t|_{\mathcal A}-\mathrm{id}_{\mathcal A}$.  Restriction to the norm unit ball is essential: an absolute bound over an unbounded algebra is defeated by scalar rescaling.  Decoherence, RG flow, and measurement back-action are different maps and must not be merged under one symbol without a model.

\paragraph{Criterion 2: Environmental Redundancy (Quantum Darwinism).}
Following quantum Darwinism \cite{Zurek2009}, let a model first supply a system/environment split, disjoint fragments $E_k$, and a classical pointer label $X$.  Redundancy can then be tested by
\begin{equation}
\chi(X:E_k)\ge(1-\delta)H(X)
\qquad\text{for many disjoint }E_k,
\end{equation}
with a finite-resolution label in continuous-spectrum cases.  Approximate commutation with an environmental algebra is not equivalent to Darwinist redundancy.  More importantly, the fragmentation that defines $E_k$ is precisely a mereological input that HAFF ultimately seeks to explain.  No factorization-free construction of $\mathcal A_E$ is given here; this is an unresolved interface tension, not a solved circularity.

\paragraph{Criterion 3: Non-scrambling Subspace.}
For normalized $O\in\operatorname{Ball}(\mathcal A)$ and a declared normalized probe class $V\in\mathfrak V$, define
\begin{equation}
C_{O,V}(t):=-\omega\!\left([O(t),V]^2\right).
\end{equation}
An informative slow-scrambling diagnostic constrains a growth rate or a post-transient window, for example
\begin{equation}
\sup_{O,V} C_{O,V}(t)\le\delta
\quad\text{for }t_0\le t\le T,
\end{equation}
where $t_0>0$ lies beyond the universal short-time regime.  Equivalently, a model may bound a fitted Lyapunov growth rate.  The former condition at merely $t\ll\tau_{\mathrm{scrambling}}$ selected almost every bounded operator at sufficiently early time and is withdrawn.

These diagnostics are observer-independent once their input data are fixed, but they are not input-free.  They may select different or even disjoint operator sectors.  In particular, Darwinist pointer records favour approximately classical commuting observables, whereas holographic geometry uses nonclassical entanglement across region algebras.  No theorem shows that one algebra satisfies both roles, and the framework must allow their intersection to be trivial.

\begin{tcolorbox}[colback=gray!5!white,colframe=black!75!black,title=\textbf{Remark: The Hamiltonian as Physical Input}]
A potential objection to the structural account of accessibility is that the interaction Hamiltonian $\hat{H}_{\text{int}}$ appears to be an arbitrary input, merely displacing observer-dependence from the algebra to the choice of Hamiltonian. We clarify that this is a misunderstanding of the framework's commitments.

The Hamiltonian is not a free parameter chosen by an observer, but a \emph{physical specification of which degrees of freedom interact and how}. In this respect, it plays a role analogous to the stress-energy tensor $T_{\mu\nu}$ in general relativity: it is input data describing the causal structure of the system, not a coordinate choice or descriptive convention.

Concretely:
\begin{itemize}
\item In quantum optics, $\hat{H}_{\text{int}}$ describes atom-photon coupling strengths and selection rules, determined by atomic energy levels and field modes.
\item In condensed matter systems, it encodes lattice structure, tunneling amplitudes, and interaction potentials, all fixed by material properties.
\item In holographic models (AdS/CFT), the boundary Hamiltonian is determined by the conformal field theory's operator content and coupling constants.
\end{itemize}

The claim is not that observers are irrelevant, but that they do not \emph{create} the supplied interaction structure.  Given a specified dynamics, norm, time window, and candidate family, whether a candidate passes a diagnostic is an objective fact.  The choice of that input family remains a separate problem.

This is not an ordinary coordinate transformation.  Different selected algebras can yield different effective operator descriptions, while the pass/fail result for a declared stability test is interaction-dependent rather than belief-dependent.

We emphasize: the framework does not solve why a particular Hamiltonian or mereology describes our universe.  It analyzes diagnostics relative to those inputs; a geometric conclusion additionally requires a reconstruction model.
\end{tcolorbox}

\paragraph{Pre-geometric Interaction Structure.}
A potential concern is that the Hamiltonian $\hat{H}_{\text{int}}$ itself presupposes spacetime structure (e.g., via spatial locality in $\hat{H} = \sum_{i,j} J_{ij} \hat{\sigma}_i \cdot \hat{\sigma}_j$), creating a circular dependence: geometry emerges from algebras determined by a Hamiltonian that already assumes geometry.

The Hamiltonian input may be \emph{metric-free} in the limited sense that it is specified as an abstract interaction graph or tensor network.  However, the nodes, tensor factors/degrees of freedom, and edge incidence are already mereological structure.  The graph avoids assuming a background metric but does not derive its own relata.  This residual input must be supplied independently and is a central quantum-mereology problem \cite{Zanardi2004,Cotler2019,CarrollSingh2021}.

Crucially, the \emph{effective spacetime geometry} that emerges from stable algebras may differ from the graph structure of the input Hamiltonian. For instance:
\begin{itemize}
\item In tensor network models (MERA), the input is a discrete causal network, but the emergent geometry can be continuous AdS space \cite{Swingle2012}.
\item In spin chain models, the Hamiltonian is defined on a 1D lattice, but entanglement structure can induce higher-dimensional effective geometry.
\item In holographic duality (AdS/CFT), the boundary Hamiltonian is defined on a fixed $(d-1)$-dimensional manifold, but the bulk geometry (including its dimensionality) emerges dynamically.
\end{itemize}

Thus, the input interaction structure can constrain candidate algebras but does not uniquely determine one.  The framework asks which candidates pass declared tests and, when a reconstruction dictionary is available, what geometric data follow.

This perspective aligns with recent work on "locality from entanglement" \cite{VanRaamsdonk2010}, where spatial locality is itself understood as arising from entanglement patterns rather than being presupposed.

\section{Relations to Existing Interpretations}
\label{sec:relations}

We now situate HAFF relative to four representative frameworks, clarifying conceptual boundaries and identifying points of potential complementarity.

\subsection{Algebraic Quantum Field Theory (AQFT)}

The closest structural affinity of HAFF is with algebraic quantum field theory \cite{Haag1996,Araki1999}. In AQFT, observable algebras are treated as primary, with states defined as positive linear functionals over these algebras. Crucially, local algebras are assigned to spacetime regions without relying on a global tensor product structure.

HAFF extends this algebraic perspective by emphasizing \emph{coarse-graining relations} between algebras. In standard AQFT, locality is typically presupposed: algebras are indexed by spacetime regions, and the split property ensures independence of spacelike-separated algebras. In HAFF, we relax this assumption and instead treat \emph{stability under physical interactions} as the criterion for algebra selection.

The transformation can be summarized as:
\begin{center}
\textbf{AQFT}: Spacetime regions $\to$ Local algebras $\to$ States \\
\textbf{HAFF}: Interaction structure $\to$ Stable algebras $\to$ Effective geometry
\end{center}

This is not a replacement of AQFT but an exploration of its structure in contexts where spacetime locality is not presupposed. The framework may be understood as asking: what happens to the algebraic approach when we do not assume a background spacetime to index our algebras?

\subsection{Relational Quantum Mechanics (RQM)}

Relational quantum mechanics \cite{Rovelli1996,LaudisaRovelli2021} emphasizes that the values of physical quantities are defined only relative to observer systems, rejecting the notion of absolute, observer-independent observables. In Rovelli's formulation, quantum mechanics is fundamentally a theory of \emph{interactions} rather than systems: what exists are relational facts, not intrinsic properties.

There is significant conceptual overlap with HAFF: both frameworks reject privileged subsystem decompositions and treat quantum descriptions as contextual. However, a key difference concerns the \emph{stabilization of relata}.

RQM analyzes relations between systems whose existence is typically taken as given (or at least presupposed operationally through interaction records). HAFF provides a mechanism for the \emph{stabilization of distinct relata} from the underlying quantum field, identifying which subsystem partitions are robustly maintained under decoherence and measurement-induced dynamics.

In this sense, HAFF may provide the stable nodes required for RQM's relational network: before relations can exist, there must be relata stable enough to participate in interactions. HAFF addresses how such stable relata emerge from the pre-factorized quantum substrate.

These are complementary rather than competing perspectives: RQM asks what observables mean relative to a system, HAFF asks which systems stabilize as distinct relata in the first place. The two frameworks operate at different levels of analysis and could in principle be combined, with HAFF providing the stability conditions under which RQM's relational structure becomes well-defined.

\subsection{QBism}

QBism \cite{Fuchs2014,FuchsMerminSchack2014} interprets quantum states as expressions of an agent's personal beliefs about measurement outcomes, emphasizing the subjective, agent-centric nature of quantum probability assignments.

Here, the distinction from HAFF is sharpest. While both frameworks reject naive realism about the quantum state, they differ fundamentally in their treatment of dependence:

\begin{itemize}
\item \textbf{QBism}: Quantum states represent personal beliefs. Dependence is epistemic and agent-centric.
\item \textbf{HAFF}: Observable algebras are selected by physical interaction structure. Dependence is structural and interaction-centric.
\end{itemize}

The key point is that interactions are not agents: they have no beliefs, make no decisions, and exist independently of any epistemic perspective.  Relative to supplied subsystem and probe data, the Hamiltonian can constrain candidate algebras; no theorem here says that it determines a unique one.

We emphasize that this is a categorical difference, not a matter of one framework being "more correct" than the other. QBism and HAFF address different questions and operate within different conceptual frameworks. The point of comparison is simply to clarify that HAFF's notion of accessibility does not reduce to QBist agent-dependence.

\subsection{Many-Worlds Interpretation (MWI)}

The Many-Worlds interpretation \cite{Wallace2012} explains the emergence of classical behavior through decoherence-induced branching, all occurring within globally unitary quantum evolution.

HAFF shares with MWI a commitment to:
\begin{itemize}
\item A single, objective global quantum state
\item Unitary evolution without collapse
\item Effective classicality emerging from entanglement structure
\end{itemize}

However, MWI typically presupposes a tensor factorization into system and environment as input, analyzing how this decomposition gives rise to branch structure. Recent work in the MWI tradition (e.g., Wallace, Saunders) addresses emergent decoherence structure, but typically within a framework where tensor factorization is assumed at the fundamental level.

HAFF complements MWI by examining the preconditions for any branching structure: before branches can emerge, there must be a notion of subsystems relative to which branching occurs. In this sense, HAFF may provide a structural framework relevant to understanding how the effective subsystem decompositions presupposed by branching emerge in the first place.

This is a point of potential contact rather than a hierarchical claim: we do not argue that MWI requires HAFF, but that the two frameworks address distinct but related aspects of quantum structure.

\section{Geometry as a Stable Organizational Phase}
\label{sec:geometry_phase}

\paragraph{Conceptual Disclaimer.}
Before discussing the analogy with condensed matter phases, we stress that it is conceptual rather than rigorous.  In condensed matter, a Hamiltonian together with thermodynamic limit, control parameters, and boundary conditions can support distinct phases and order parameters.  HAFF currently lacks an analogous proved selection model.  The comparison is therefore a heuristic for stability and contingent objectivity, not an independent result.

\paragraph{Contingent Objectivity.}
A better analogy is spontaneous symmetry breaking.  Below the Curie temperature, one ferromagnetic Hamiltonian admits macroscopically distinct phases with magnetization $+M$ or $-M$.  Boundary conditions or a small external field select a phase; once selected, its magnetization is objective and reproducible, although the microscopic laws do not choose a unique sign in isolation.

Likewise, a supplied interaction/mereology datum may admit several stable candidate algebras.  A physical selection condition can make one candidate salient without making it observer-created.  If a separate reconstruction dictionary maps that candidate's entanglement data to a metric, observers using the same inputs obtain the same reconstructed result.  The analogy supports contingent objectivity only; it neither proves algebra selection nor supplies the reconstruction dictionary.

This makes the intended interpretation coherent without proving it: reconstructed geometry may be contingent yet objective within a model, much as an order parameter is objective in a selected phase.

\subsection{The Phase Analogy}

With these caveats in place, we develop the analogy with condensed matter phases.

In statistical mechanics, different thermodynamic phases (solid, liquid, gas, magnetic, superconducting) represent distinct organizational patterns of microscopic degrees of freedom. These phases are:
\begin{itemize}
\item \textbf{Emergent}: They arise from collective behavior, not from individual constituents.
\item \textbf{Stable}: They persist under perturbations below characteristic energy scales.
\item \textbf{Observable}: They manifest in measurable order parameters.
\item \textbf{Contingent}: They depend on external conditions (temperature, pressure, fields).
\end{itemize}

We propose that emergent geometry in HAFF exhibits analogous features:

\begin{table}[h]
\centering
\begin{tabular}{|l|l|}
\hline
\textbf{Condensed Matter} & \textbf{HAFF} \\
\hline
    Hamiltonian $\hat{H}$ & Metric-free but mereologically supplied interaction graph \\
    Control parameter (Temperature $T$) & Model-specific coupling/noise parameter \\
Symmetry breaking & Algebra selection \\
Order parameter $\langle M \rangle$ & Entanglement pattern $I(A:B)$ \\
Phase transition & Geometry emergence \\
    Critical temperature $T_c$ & Model-specific stability threshold \\
\hline
\end{tabular}
\caption{Heuristic analogy between condensed matter phases and candidate algebra selection. No universal HAFF control parameter or critical threshold is derived.}
\label{tab:phase}
\end{table}

In holographic models (AdS/CFT), a model-specific analogue is the ratio $\ell_{\text{AdS}}/\ell_{\text{Planck}}$, which controls the validity of a classical bulk description \cite{Maldacena1999}; this does not supply a universal HAFF control parameter.

Just as ferromagnetism is a stable organizational pattern of spin alignment below the Curie temperature, geometry may be modelled as a stable organizational pattern only when explicit algebra-selection and reconstruction conditions are supplied.

\paragraph{Post-selection entanglement diagnostic.}
After a model has supplied subsystems and an interaction graph $G=(V,E)$, one may define the average pair mutual information
\begin{equation}
\rho_{\mathrm{ent}}^{(G)}
:=\frac{1}{|E|}\sum_{(i,j)\in E} I(i:j),
\end{equation}
for $|E|>0$.  If all unordered pairs are used, the denominator is $\binom{|V|}{2}$ instead.  The earlier normalization by $N$ was undefined without a graph-scaling convention.

This quantity depends on the very subsystem graph whose origin is at issue, so it cannot select the pre-geometric mereology without circularity.  It is only a diagnostic after that input is fixed.  Low average pair mutual information does not by itself imply an area law or smooth geometry, and high values are not equivalent to a scrambling rate.

In a particular holographic model, quantities such as $\ell_{\mathrm{AdS}}/\ell_P$ control a semiclassical expansion \cite{Maldacena1999}; they are not identified here with $\rho_{\mathrm{ent}}^{(G)}$.

Whether any model exhibits a genuine transition in geometric reconstructibility is an open, model-specific question.

\subsection{Geometric Admissibility}

Not all coarse-graining structures admit geometric interpretation. As discussed in our previous work (Definition 4.2), we require geometric admissibility conditions:
\begin{enumerate}
\item Finite correlation length (exponentially decaying correlations)
\item Monotonic decay of mutual information under refinement
\item Stability under perturbations
\end{enumerate}

These conditions are not arbitrary but reflect empirical observations from known emergent geometries:
\begin{itemize}
\item \textbf{AdS/CFT}: Boundary states with area-law entanglement give rise to smooth bulk geometries \cite{RyuTakayanagi2006}.
\item \textbf{Tensor networks}: MERA and similar structures with finite bond dimension naturally induce geometric connectivity \cite{Swingle2012}.
\item \textbf{Condensed matter}: Ground states of local Hamiltonians typically satisfy area laws and admit geometric descriptions.
\end{itemize}

Geometry, in this view, is conjectured to represent a special organizational phase characterized by specific entanglement structure.  Demonstrating an attractor requires a concrete coarse-graining dynamics and reconstruction dictionary and remains open.

\section{Scope and Future Directions}
\label{sec:scope}

\subsection{Structural Assumptions}

The framework rests on three core assumptions:
\begin{enumerate}
\item \textbf{Global state objectivity}: There exists a universal quantum state $|\Psi_U\rangle$, independent of observers.
\item \textbf{Algebraic priority}: Observable-algebra data are treated as prior to the effective tensor factorizations reconstructed from them; their own physical selection remains open.
\item \textbf{Stability-based emergence}: Stable candidate algebras are proposed as inputs to effective subsystem and geometric reconstruction, conditional on the missing selection and dictionary results.
\end{enumerate}

\subsection{What the Framework Does NOT Commit To}

To prevent interpretational overreach, we explicitly enumerate what the framework does \emph{not} claim:

\begin{itemize}
\item That observers create reality or that consciousness plays a fundamental role
\item That spacetime is illusory or that geometry has no objective existence
\item That quantum mechanics is incomplete or requires modification
\item That the framework solves the measurement problem
\item That Buddhist metaphysics or any other philosophical tradition is presupposed
\end{itemize}

The framework is structurally neutral regarding these questions. It analyzes consequences of relaxing the assumption of canonical tensor factorization, but does not commit to any particular metaphysical position beyond what the formalism requires.

\subsection{Open Questions}

The following questions include load-bearing open gaps as well as later empirical directions.  The mereology, algebra-selection, and reconstruction problems must be solved in concrete models before the heuristic diagnostics become a derivation.

\begin{enumerate}
\item \textbf{Quantum mereology}: What physical datum selects the nodes/factors or candidate algebra family without presupposing the subsystem structure to be explained \cite{Zanardi2004,Cotler2019,CarrollSingh2021}?

\item \textbf{Dynamical algebra selection}: Can interaction Hamiltonians, coupling strengths, or network topologies be shown to select particular stable algebras over others?

\item \textbf{Information-theoretic criteria}: Do preferred geometric descriptions correlate with minimal description length, robustness under noise, or computational accessibility?

\item \textbf{Quantum field theory extensions}: How does the framework interact with locality structures in algebraic QFT? Can continuum limits be rigorously constructed?

\item \textbf{Empirical signatures}: Do different coarse-graining choices lead to distinguishable effective descriptions in semiclassical regimes or quantum gravity-motivated models?

\item \textbf{Connections to quantum complexity}: How does the framework relate to recent work on complexity-based approaches to spacetime emergence?
\end{enumerate}

These questions are intentionally left open. The present work provides a structural scaffold within which they can be formulated precisely, but does not claim to resolve them.

\subsection{Beyond Structural Analysis}

The framework developed here is deliberately limited to structural and interpretational analysis. Broader philosophical implications—concerning causation, free will, ontology, and connections to contemplative traditions—lie beyond the scope of this technical paper. Such questions are addressed in a companion philosophical essay currently in preparation.

\section{Conclusion}

We have clarified the conceptual commitments of the Holographic Alaya-Field Framework, addressing three potential misreadings:

\begin{enumerate}
\item \textbf{Against subjectivism}: Accessibility is defined structurally via stability conditions determined by physical interaction structure, not by observer beliefs or epistemic states.

\item \textbf{Against anti-realism}: Emergent geometry is a stable, measurable feature of entanglement structure—contingent on physical conditions but objective within those conditions, analogous to thermodynamic phases.

\item \textbf{Against trivialism}: Coarse-graining dependence reflects genuine physical structure (accessible algebras), not mere coordinate choice. Different algebras can induce inequivalent geometries when a supplied reconstruction dictionary separates the compared profiles modulo diffeomorphisms.
\end{enumerate}

The framework has been situated relative to AQFT (closest affinity), RQM (complementary), QBism (categorically distinct), and MWI (potentially complementary). Throughout, we have emphasized what the framework commits to structurally and what it leaves open interpretionally.

The central insight is modest but consequential: by removing tensor factorization from fundamental assumptions and treating it as emergent from coarse-graining structure, we reveal that geometry is more context-dependent than typically acknowledged—not in an epistemic or observer-relative sense, but in a structural, interaction-dependent sense.

This perspective does not resolve deep problems in quantum foundations or quantum gravity, but it clarifies the conditions under which subsystem structure and geometry emerge. By making explicit an assumption that is often left implicit, we hope to have opened new avenues for investigating the relationship between quantum states, observable structure, and spacetime.

\section*{Acknowledgments}

The author thanks the anonymous reviewers for their insightful comments and suggestions, which greatly improved the clarity and rigor of this work.

\begin{thebibliography}{99}

\bibitem{Liu2026}
S.~Liu, \emph{Emergent Geometry from Coarse-Grained Observable Algebras: The Holographic Alaya-Field Framework}, Zenodo (2026), DOI: 10.5281/zenodo.18361706.

\bibitem{Zurek2009}
W. H. Zurek, \emph{Quantum Darwinism}, Nature Physics \textbf{5}, 181 (2009).

\bibitem{Hayden2007}
P. Hayden and J. Preskill, \emph{Black holes as mirrors: quantum information in random subsystems}, JHEP \textbf{09}, 120 (2007).

\bibitem{VanRaamsdonk2010}
M. Van Raamsdonk, \emph{Building up spacetime with quantum entanglement}, Gen. Relativ. Gravit. \textbf{42}, 2323 (2010).

\bibitem{Swingle2012}
B. Swingle, \emph{Entanglement Renormalization and Holography}, Phys. Rev. D \textbf{86}, 065007 (2012).

\bibitem{Haag1996}
R. Haag, \emph{Local Quantum Physics: Fields, Particles, Algebras}, Springer-Verlag (1996).

\bibitem{Araki1999}
H. Araki, \emph{Mathematical Theory of Quantum Fields}, Oxford University Press (1999).

\bibitem{Rovelli1996}
C. Rovelli, \emph{Relational Quantum Mechanics}, Int. J. Theor. Phys. \textbf{35}, 1637 (1996).

\bibitem{LaudisaRovelli2021}
F. Laudisa and C. Rovelli, \emph{Relational Quantum Mechanics}, Stanford Encyclopedia of Philosophy (2021).

\bibitem{Fuchs2014}
C. A. Fuchs, \emph{QBism, the Perimeter of Quantum Bayesianism}, arXiv:1003.5209 (2014).

\bibitem{FuchsMerminSchack2014}
C. A. Fuchs, N. D. Mermin, and R. Schack, \emph{An Introduction to QBism with an Application to the Locality of Quantum Mechanics}, Am. J. Phys. \textbf{82}, 749 (2014).

\bibitem{Wallace2012}
D. Wallace, \emph{The Emergent Multiverse}, Oxford University Press (2012).

\bibitem{RyuTakayanagi2006}
S. Ryu and T. Takayanagi, \emph{Holographic Derivation of Entanglement Entropy from AdS/CFT}, Phys. Rev. Lett. \textbf{96}, 181602 (2006).

\bibitem{Maldacena1999}
J. M. Maldacena, \emph{The Large N limit of superconformal field theories and supergravity}, Int. J. Theor. Phys. \textbf{38}, 1113 (1999).

\bibitem{Zanardi2004}
P.~Zanardi, D.~A.~Lidar, and S.~Lloyd,
\emph{Quantum tensor product structures are observable induced},
Phys.\ Rev.\ Lett.\ \textbf{92}, 060402 (2004).

\bibitem{Cotler2019}
J.~Cotler, G.~Penington, and D.~H.~Ranard,
\emph{Locality from the spectrum},
Commun.\ Math.\ Phys.\ \textbf{368}, 1267--1296 (2019).

\bibitem{CarrollSingh2021}
S.~M.~Carroll and A.~Singh,
\emph{Quantum mereology: Factorizing Hilbert space into subsystems with quasiclassical dynamics},
Phys.\ Rev.\ A \textbf{103}, 022213 (2021).

\end{thebibliography}

\end{document}
